\documentclass[11pt]{article}

\usepackage[margin=1in]{geometry}
\usepackage{hyperref}
\usepackage{amsmath}
\usepackage{enumitem}
\usepackage{listings}
\usepackage{xcolor}

\hypersetup{
    colorlinks=true,
    linkcolor=blue,
    urlcolor=blue,
    citecolor=blue
}

\title{CS 300 -- Week 2 Resources\\\large Sorting Algorithms and Algorithm Analysis}
\author{Jose de Lima}
\date{\today}

\begin{document}
\maketitle

\section*{Overview}
Week 2 focuses on two foundational comparison-based sorting algorithms --
\textbf{selection sort} and \textbf{insertion sort} -- and on the language we
use to describe how algorithms scale: \textbf{Big-O notation}. The resources
below introduce each algorithm step-by-step and provide a reference for the
asymptotic running times of the most common data-structure operations.

\section{Selection Sort}
\textbf{Resource:} \href{https://www.youtube.com/watch?v=g-PGLbMth_g}{Selection Sort in 3 Minutes (YouTube)}

\subsection*{Idea}
Selection sort builds the sorted portion of the array one element at a time by
repeatedly \emph{selecting} the smallest (or largest) element from the
unsorted portion and swapping it into place.

\subsection*{Algorithm}
\begin{enumerate}[leftmargin=*]
    \item Treat index $0$ as the boundary between the sorted and unsorted halves.
    \item Scan the unsorted half to find the index of the minimum element.
    \item Swap that minimum with the element at the boundary.
    \item Advance the boundary by one and repeat until the array is sorted.
\end{enumerate}

\subsection*{Pseudocode}
\begin{lstlisting}[basicstyle=\ttfamily\small,frame=single]
for i = 0 to n - 2:
    minIndex = i
    for j = i + 1 to n - 1:
        if A[j] < A[minIndex]:
            minIndex = j
    swap(A[i], A[minIndex])
\end{lstlisting}

\subsection*{Complexity}
\begin{itemize}[leftmargin=*]
    \item \textbf{Time:} $O(n^2)$ in the best, average, and worst cases, because
          the inner scan always traverses the remaining unsorted region.
    \item \textbf{Space:} $O(1)$ -- sorting is done in place.
    \item \textbf{Stability:} Not stable by default (equal keys can be reordered
          by the swap step).
    \item \textbf{Swaps:} At most $n-1$, which is useful when writes are
          expensive (e.g., flash memory).
\end{itemize}

\section{Insertion Sort}
\textbf{Resource:} \href{https://www.youtube.com/watch?v=JU767SDMDvA}{Insertion Sort in 2 Minutes (YouTube)}

\subsection*{Idea}
Insertion sort grows a sorted prefix by taking the next unsorted element and
\emph{inserting} it into its correct position among the already-sorted items,
shifting larger elements one slot to the right to make room.

\subsection*{Algorithm}
\begin{enumerate}[leftmargin=*]
    \item Start at index $1$; treat index $0$ as a trivially sorted prefix.
    \item Save $A[i]$ as \texttt{key}.
    \item Walk backwards through the sorted prefix, shifting any element
          greater than \texttt{key} one position to the right.
    \item Place \texttt{key} in the resulting gap.
    \item Advance $i$ and repeat until the array is sorted.
\end{enumerate}

\subsection*{Pseudocode}
\begin{lstlisting}[basicstyle=\ttfamily\small,frame=single]
for i = 1 to n - 1:
    key = A[i]
    j = i - 1
    while j >= 0 and A[j] > key:
        A[j + 1] = A[j]
        j = j - 1
    A[j + 1] = key
\end{lstlisting}

\subsection*{Complexity}
\begin{itemize}[leftmargin=*]
    \item \textbf{Best case:} $O(n)$ when the input is already (nearly) sorted
          -- the inner loop exits immediately each iteration.
    \item \textbf{Average / worst case:} $O(n^2)$ for random or reverse-sorted
          input.
    \item \textbf{Space:} $O(1)$ -- in-place.
    \item \textbf{Stability:} Stable -- equal keys keep their original order.
    \item \textbf{Practical note:} Very fast on small or nearly-sorted arrays,
          which is why hybrid sorts (e.g., Timsort, Introsort) fall back to
          insertion sort for small partitions.
\end{itemize}

\section{Selection Sort vs. Insertion Sort}
\begin{center}
\begin{tabular}{l|l|l}
\textbf{Property} & \textbf{Selection Sort} & \textbf{Insertion Sort} \\ \hline
Best-case time    & $O(n^2)$                & $O(n)$                  \\
Average time      & $O(n^2)$                & $O(n^2)$                \\
Worst-case time   & $O(n^2)$                & $O(n^2)$                \\
Space             & $O(1)$                  & $O(1)$                  \\
Stable?           & No                      & Yes                     \\
\# of swaps       & $O(n)$                  & $O(n^2)$                \\
Adaptive?         & No                      & Yes                     \\
\end{tabular}
\end{center}

\section{Big-O Cheatsheet}

\subsection*{What Big-O Describes}
Big-O notation expresses an \emph{upper bound} on the growth rate of an
algorithm's running time (or space) as a function of the input size $n$. It
lets us compare algorithms independent of hardware, language, or constant
factors.

\subsection*{Common Growth Classes (fastest to slowest)}
\begin{itemize}[leftmargin=*]
    \item $O(1)$ -- constant: hash-table lookup, array index access.
    \item $O(\log n)$ -- logarithmic: binary search, balanced-BST operations.
    \item $O(n)$ -- linear: a single pass over $n$ elements.
    \item $O(n \log n)$ -- linearithmic: merge sort, heap sort, quicksort (avg).
    \item $O(n^2)$ -- quadratic: selection sort, insertion sort, bubble sort.
    \item $O(2^n)$ -- exponential: naive recursion over subsets.
    \item $O(n!)$ -- factorial: brute-force permutations (e.g., TSP).
\end{itemize}

\subsection*{Related Notation}
\begin{itemize}[leftmargin=*]
    \item $\Omega(f(n))$ -- asymptotic lower bound (best case).
    \item $\Theta(f(n))$ -- tight bound (matching upper and lower bound).
\end{itemize}

\subsection*{Rules of Thumb}
\begin{itemize}[leftmargin=*]
    \item Drop constants: $O(3n) = O(n)$.
    \item Drop non-dominant terms: $O(n^2 + n) = O(n^2)$.
    \item Nested loops over the same collection usually multiply: two nested
          loops over $n$ items yield $O(n^2)$.
    \item Halving the problem each step usually yields $O(\log n)$.
\end{itemize}

\subsection*{Why it Matters for This Course}
When we compare vectors, linked lists, hash tables, and trees later in CS 300,
the deciding factor is rarely which data structure \emph{can} store the data
-- it is which one makes the operations we care about (insert, search, remove,
traverse) cheap as $n$ grows. The cheatsheet summarizes those trade-offs in a
single table and should be kept close at hand throughout the term.

\section*{Links Summary}
\begin{itemize}[leftmargin=*]
    \item Selection Sort video: \url{https://www.youtube.com/watch?v=g-PGLbMth_g}
    \item Insertion Sort video: \url{https://www.youtube.com/watch?v=JU767SDMDvA}
\end{itemize}

\end{document}
